% Paper 3 — Introduction Section
% Building Local-Language Economic Narrative Indices

\section{Introduction}

Text-as-data methods have transformed empirical research in economics, political science, and sociology \citep{gentzkow2019text, grimmer2019text}. Researchers now routinely extract sentiment, topics, and narrative frames from news articles, social media, and policy documents to measure economic conditions, predict financial markets, and understand public discourse \citep{tetlock2007sad, baker2016measuring}. Yet this transformation has been overwhelmingly English-language: 84\% of published studies analyze English text, and 56\% focus exclusively on the United States \citep{nabil2026economic}.

This linguistic concentration creates a fundamental measurement gap. For the 5.5 billion people who do not speak English as their primary language, local-language economic narratives remain only weakly visible to researchers and policymakers. In Bangladesh---a large developing economy and the centre of a 265-million-speaker Bangla language community---no reproducible economic narrative index has been available. Official macroeconomic series are released with delay, informal activity is only partially measured, and financial-market signals are thinner than in advanced economies. Yet Bangla newspapers publish daily accounts of prices, exchange rates, trade, fiscal policy, labour markets, banking stress, and household economic pressure.

This paper addresses a simple question: \textit{can local-language news be transformed into a validated economic narrative index, and what is the minimum annotation investment required to do so reliably?}

We introduce the BENI v1 method pipeline, a reproducible workflow for building a Bangla Economic Narrative Index from harmonised Bangla news data. The BENI v1 database contains 1,467,705 article records assembled from dated Potrika articles from 2014--2020 and post-2020 BNAD articles, with standardised metadata, deduplication fields, text hashes, economic seed labels, and model-ready text. The frozen validation set used in this paper contains 186 clean reference labels mapped to canonical BENI v1 article IDs, while 114 ambiguous rows are excluded from clean evaluation.

The current frozen BENI v1 prediction run produces article-level economic relevance scores for the repaired corpus and supports two monthly index variants: article-weighted and source-balanced. The source-balanced specification is the default measurement series, because it reduces domination by high-volume outlets.

Our contribution is threefold. First, we document the first reproducible BENI v1 database and measurement pipeline for Bangla economic narratives. Second, we provide a frozen validation benchmark for Bangla economic text classification: TF-IDF + Logistic Regression achieves 88.17\% accuracy and 0.823 F1 on the Economic class against the 186-row clean validation set. Third, we show how a harmonised low-resource corpus can be turned into a frozen article-level prediction file and a monthly narrative index that can be reused in later macroeconomic validation and nowcasting papers.

The paper proceeds as follows. Section~\ref{sec:related} reviews related work on text-as-data in economics and local-language NLP. Section~\ref{sec:data} describes the BENI v1 database and annotation protocol. Section~\ref{sec:methods} presents the classification, active-learning, index-construction, and validation pipeline. Section~\ref{sec:results} reports calibration results. Section~\ref{sec:discussion} discusses implications, and Section~\ref{sec:template} provides a replicable template for other low-resource languages.
